\documentclass[11pt]{article}
\usepackage[margin=1in]{geometry}
\usepackage{amsmath,amssymb,amsthm}
\usepackage{booktabs,array,enumitem}
\usepackage[hidelinks]{hyperref}
\setlist{itemsep=2pt,topsep=3pt}
\newcommand{\st}[1]{\textsf{\small[#1]}}
\title{Erd\H{o}s Problem \#644: progress report}
\author{Project status as of 24 November 2026}
\date{}
\begin{document}
\maketitle

\begin{abstract}
We work on the question whether $f(k,7)=(\tfrac34+o(1))k$, where $f(k,7)$ is the largest transversal number of a
$k$-uniform family in which every seven members can be pierced by two points. The known bounds are
$\tfrac34k+1\le f(k,7)\le\lceil 7k/8\rceil$ (Fon-Der-Flaass--Kostochka--Woodall). \textbf{We have not solved the
problem.} We have a verified reduction of it to a covering game confined to at most $3k$ points. We have closed the
crowded and spread regions of that game, and proved that every \emph{structured} counterexample is impossible in the
remaining region. What is left is a precise statement about ``random-looking'' families in a window of
$1.75k$--$2.25k$ points. This report says which parts are publishable and which are only working evidence.
\end{abstract}

\section*{Status labels}
\st{LEAN}: machine-checked in Lean~4/Mathlib, standard axioms only.
\st{VERIFIED}: a hand proof that an independent agent, given only the statement and the proof, tried to break.
That agent checked every step and ran its own adversarial exact tests.
\st{EXACT}: rational or brute-force computation on finite instances.
\st{FLOAT}: floating-point evidence.
\st{DEAD}: refuted or shown not to apply.
The status of every claim is tracked in \texttt{notes/LEDGER.md}.

\section{The approach in one page}
\textbf{Request principle.} If $\tau(H)>T$, every set $S$ of $T$ points misses some edge. A ``requester'' names sets
$S_1,S_2,\dots$, and an ``adversary'' answers each with an edge avoiding it. If the requester can always collect seven
answers with no 2-point transversal, then $f(k,7)\le T$. With unboundedly many requests the adversary is no stronger
than a fixed family (CORE \S29), so the game \emph{is} the problem, not a relaxation.

\textbf{Reduction.} Three answers with empty common intersection confine all later play to their union $V$, with
$|V|\le 3k$. Requester wins in the confined game transfer to the full game (TWIN \S13) \st{VERIFIED}. In units of $k$
we write $t=T/k=\tfrac34+\varepsilon$ and $N=|V|/k$. We pass to \emph{co-edges} $K=V\setminus B$. Seven answers have
no 2-cover iff their co-edges cover every pair of $V$, i.e.\ the sets $I_x=\{i:x\in K_i\}$ are pairwise intersecting.
And $\tau>T$ iff the co-edge family is \emph{$T$-covering}: every $T$-set lies in some co-edge.

\textbf{Case split by $N$.}
\begin{center}\small
\begin{tabular}{@{}p{4.6cm}p{5.6cm}p{4.6cm}@{}}\toprule
region & content & status\\\midrule
$N\le 3t-\tfrac12$ (crowded) & Proposition 16.A & \st{VERIFIED} at first order\\
$3t-\tfrac12<N<3-t$ (\textbf{core window}) & statement CP / $(\Sigma)$ below & \textbf{open}\\
$N\ge 3-t$ (petal side) & petal reduction, Theorems J, K, L, AP, PS--PS3 & mostly \st{VERIFIED}; two residues\\\bottomrule
\end{tabular}
\end{center}

\section{What is established}
\subsection{Machine-checked}
About twenty lemmas are in Lean with no \texttt{sorry}. They cover:
\begin{itemize}
\item the tight-triple lemma, a transversal bound we call Theorem~A, stability, the lopsided-top bound, the
sparse-triple lemma, the trace lemma, the pair gap and the three-centres lemma;
\item the Fano criterion and Fano blocks, the flat 7-tuple lemma (top-degree requests give seven edges with every
point of degree at most~4), and the waste adversary;
\item the anchor and petal reductions;
\item Lemmas 4.1--4.8 of the unrefereed $6/7$ manuscript.
\end{itemize}
The headline statements are frozen in \texttt{Statements.lean}. The main $6/7$ theorem itself is \emph{not}
formalised.

\subsection{Hand proofs that passed independent verification}
\begin{itemize}
\item \textbf{Tools of the core window:} the slack count, Theorem T (triangle finish), Theorem Q, which is the same as
its set form Q$'$, and the reduction CP~$\Leftarrow(\Sigma)$ with one clause corrected by the verifier.
\item \textbf{Classification of structured families (core window):}
 \begin{itemize}
 \item \emph{Theorem S}: no family whose membership depends only on profiles of a bounded partition is a
 counterexample.
 \item \emph{Theorem S$^+$}: the same at finite $k$, for up to $4\varepsilon k-12$ cells of any size.
 \item \emph{\S46}: ordered frozen rules lose.
 \item \emph{Theorem W}: any greedy priority-order adversary whose choice of order depends on the request only through
 a window of $\le T-k/2$ points loses in seven requests, with any number of orders.
 \end{itemize}
\item \textbf{Petal side:} Theorems J, J$'$, J$''$, K, L, N, DS-b, AP, PS, PS2, PS3 and Lemma FF.
\end{itemize}
The verifiers ran several hundred thousand exact or brute-force games, including adversaries built against the
weakest step of each proof. They found no counterexample. They did find corrections:
\begin{itemize}
\item The pool theorem P and Proposition 16.A hold at first order but need repair at finite $k$. P fails at
$(k,T,N)=(5,4,9)$ because $C(9,4,2)=8$.
\item A short count (CORE 21.2) that was meant to close the bottom strip $3t-\tfrac12<N\le 5t-2$ has a gap. That strip
is therefore part of the open window.
\item Double sealing (DS) fails when the adversary may leak outside its arenas. This was repaired by PS--PS3.
\end{itemize}

\subsection{Exact small values}
With 8 requests, every 6- and 7-uniform family with $\tau=k$ contains seven members with no 2-cover
\st{EXACT certificate, independent checker}. Together with TWIN \S13, now verified, this gives $f(6,7)=5$ and
$f(7,7)=6$, against the general upper bounds 6 and 7. The matching lower bounds rest on covering numbers $C(v,k,2)$
from our own ILP. They still have to be checked against the La~Jolla tables.

\section{What is publishable, and what is not}
\textbf{Worth publishing or presenting}, in our judgement:
\begin{enumerate}
\item \emph{The reduction and game equivalence.} The problem equals a covering game confined to $\le3k$ points, and
bounded-depth games are exact truncations of it \st{VERIFIED}. This is a clean reformulation that others can use.
\item \emph{The classification theorems S, S$^+$, \S46, W.} These are genuine theorems about the problem: a
counterexample to $3/4$ cannot be a bounded blow-up, cannot follow fixed orders, and cannot choose among fixed orders
by a small window. They constrain any future construction. \st{VERIFIED}
\item \emph{The exact statement of what remains}, $(\Sigma)$, together with the proved negative results:
 \begin{itemize}
 \item junta theorems (Dinur--Friedgut, Keller--Lifshitz) do not apply, because the families are far too sparse;
 \item spread and sunflower lemmas do not apply, because the families are at most $2.3$-spread;
 \item exact symmetrisation destroys the property one needs.
 \end{itemize}
 This saves the next person months.
\item \emph{Lean formalisation} of the FKW-type lemmas and of the $6/7$ manuscript's lemmas. It is valuable as
verification of others' unrefereed work.
\item \emph{$f(6,7)=5$ and $f(7,7)=6$}, after the covering-number check: small, but new as exact values as far as we
know \st{cite?}.
\end{enumerate}
\textbf{Not worth publishing on its own:}
\begin{itemize}
\item The numerical game searches: the tilted-spread strategy trees, the depth-2/3 values and the SPOIL experiments.
They are float-level, narrow, or superseded by the theorems above.
\item The small-$k$ adversary wins (e.g.\ $(10,8)$): we showed these are integrality artefacts.
\item The parked ``$0.84$'' upper-bound candidate: only partly certified, and weaker than the claimed $6/7$.
\item The petal theorems, as a separate paper: they are correct, but specialised and meaningful only as part of a
full solution.
\end{itemize}
\textbf{Honest headline.} There is no new upper bound on $f(k,7)$ and no solution. The publishable content is a
verified reformulation plus impossibility results for structured counterexamples.

\section{What remains}
\begin{enumerate}
\item \textbf{Core window (the main obstacle).} We must show that every $T$-covering co-edge family on $N$ points,
$3t-\tfrac12<N<3-t$, contains seven members covering all pairs. By the reduction it suffices to find one of three
shapes: a Q$'$ quadruple, a pencil triple with union $V$ and large triple intersection, or a Theorem-T triangle. The
families not yet handled are:
 \begin{itemize}
 \item hash-like ones, whose choice reads more than about $k/4$ requested points;
 \item non-greedy ones;
 \item self-cut ones, whose cells are cut out by their own earlier members.
 \end{itemize}
 A decomposition in the style of Kupavskii--Zakharov reduces everything to the hash-like case.
\item \textbf{Petal strip.} The pooled arena adversary, for $\tfrac34<t<\tfrac79$ and
$\tfrac t2-\tfrac16\le\delta<1-t$. The sealing method is provably maxed out there. The remaining question is whether
an \emph{adaptive} final request saves one $\delta$ of sealing cost.
\item \textbf{Unusual petal families.} Two-class families not of arena-with-pool form. An exact search at $k=12,16$
found none.
\end{enumerate}

\section{What writing this made clear}
Writing the case split out in one place showed three things we had not stated sharply before.

\textbf{(a) A counterexample would need extreme negative correlation.} For the random threshold $T$-covering family
in the window, the expected number of good sevens is $e^{ck}$ with $c$ between $2.1$ and $3.4$, and the expected
number of pencils is also exponentially large \st{FLOAT, CORE \S57 code}. So a counterexample must avoid exponentially
many candidates that a random family would contain. Our structure theorems kill the structured ways of doing that.
Avoiding a fixed configuration in a sparse, random-looking regime is exactly what the \textbf{hypergraph container
method} (Balogh--Morris--Samotij; Saxton--Thomason) is built for. Containers need a \emph{supersaturation} statement:
a $T$-covering family that is far from every container must contain many good sevens. This is a route we have not
tried, and the first-moment numbers suggest it is the natural one.

\textbf{(b) The threshold $k/4$ appears everywhere for one reason.} Theorem W stops at windows of $T-k/2=(\tfrac14+
\varepsilon)k$ points, Theorem S$^+$ at $4\varepsilon k$ cells, and the pool families at $t=\tfrac34$ exactly. In each
case the slack is the amount by which $T$ exceeds $3k/4$. This is consistent with $3/4$ being the true constant, and
it locates the difficulty in a simple place: the open case is adversaries that use more than the slack's worth of
information about the request.

\textbf{(c) The two open pieces are alike.} The petal strip and the core window both reduce to whether an
\emph{adaptive} last move by the requester beats an adversary that hides information. Both may need the same
ingredient: an averaging or entropy argument over the requester's options, rather than a fixed plan.

\section{Work in progress}
Agents now running, each held to the same verification gate:
\begin{itemize}
\item an entropy-method attack on the core window (in the style of Gilmer's union-closed argument);
\item a flag-algebra / SDP relaxation of the $k\to\infty$ limit object;
\item a fractional relaxation followed by rounding, via $(p,q)$-type tools;
\item an exact decision on a $k=8$ family in which no good seven has been found yet;
\item the adaptive petal-strip game;
\item the structure lemma for the unusual petal families.
\end{itemize}
Suggestion (a), hypergraph containers, is new and is not yet assigned.

\end{document}
